\chapter{Rosacea}

\section*{Definition}
Rosacea is a chronic, relapsing inflammatory disorder of the facial skin characterized by transient or persistent erythema, telangiectasias, papules, pustules, and occasionally phymatous changes. It affects approximately 5--10\% of adults, predominantly fair-skinned individuals of Celtic or Northern European descent.

\section*{What Happens in Your Body}
The disease development involves dysregulation of the innate immune system with overexpression of cathelicidin (LL-37) and kallikrein 5, triggering aberrant inflammatory responses to environmental triggers. Demodex folliculorum mite overgrowth, vascular hyperreactivity, and activation of TLR2 on keratinocytes contribute to the pro-inflammatory state. Ultraviolet radiation induces reactive oxygen species and matrix metalloproteinases, worsening cutaneous inflammation.

\section*{Causes}
\begin{itemize}
  \item \textbf{Environmental triggers:} UV exposure, heat, cold, strong winds, spicy foods, alcohol, hot beverages
  \item \textbf{Biological:} \textit{Demodex folliculorum} colonization, \textit{Bacillus oleronius} (Demodex-associated), \textit{Helicobacter pylori} (controversial association)
  \item \textbf{Vascular:} Genetic predisposition to flushing, rosacea-associated HLA alleles
  \item \textbf{Endocrine:} Menopause, thyroid disease (possible association)
  \item \textbf{Medications:} Topical corticosteroids (steroid rosacea), niacin, vasodilators
\end{itemize}

\section*{When to See a Doctor}
See your doctor if:
\begin{itemize}
  \item Persistent facial erythema, frequent flushing, or papulopustular eruption
  \item Ocular symptoms: dry eye, blepharitis, conjunctival injection, corneal irritation, or vision changes
  \item Rhinophyma or other phymatous changes causing cosmetic or functional concern
  \item Worsening despite avoidance of triggers and over-the-counter therapies
\end{itemize}

\section*{Related Symptoms}
\begin{itemize}
  \item Acne vulgaris, seborrheic dermatitis, perioral dermatitis, cutaneous lupus erythematosus
\end{itemize}
\textbf{Important:} Ocular rosacea occurs in 50\% of people and may precede cutaneous manifestations; untreated, it can progress to keratitis, corneal ulceration, and vision loss.

\section*{Self-Care / Home Management}
\begin{itemize}
  \item Identify and avoid personal trigger factors; keep a symptom diary to track flares
  \item Gentle skin care: mild, non-abrasive cleansers, fragrance-free moisturizers, and daily broad-spectrum SPF 50+ sunscreen
  \item Avoid topical corticosteroids, harsh exfoliants, and astringents on facial skin
\end{itemize}

\section*{How Your Doctor Will Evaluate You}
\begin{itemize}
  \item Clinical diagnosis based on presence of one or more diagnostic features: flushing (transient erythema), non-transient erythema, papules/pustules, telangiectasias, phymatous changes
  \item Subtype classification: erythematotelangiectatic (ETR), papulopustular, phymatous, ocular
  \item Skin biopsy rarely needed; consider direct immunofluorescence if lupus erythematosus is in differential
\end{itemize}

\section*{Treatment Options}
\begin{itemize}
  \item Mild--moderate: Topical metronidazole 0.75--1\% cream/gel BID, azelaic acid 15\% gel BID, or ivermectin 1\% cream daily
  \item Inflammatory papules/pustules: Topical brimonidine 0.33\% gel or oxymetazoline 1\% cream for erythema; oral doxycycline 40 mg (anti-inflammatory dose daily) or minocycline
  \item Phymatous: Laser therapy (pulsed-dye, KTP, CO2), electrosurgery, or dermabrasion
  \item Ocular: Warm compresses, artificial tears, oral doxycycline; ophthalmology referral for corneal involvement
\end{itemize}

\clearpage
